% ECONOMICS_PROBLEM: {"id":"C71-12","title":"Grand-coalition arrival in the patient limit","jel_code":"C71","source_id":"OP-0222"}
% FIRST_POSTED: 2026-10-09
% ATTEMPT_STATUS: unresolved
% ENTRY_KIND: research_attempt
\documentclass[11pt]{article}
\usepackage{amsmath,amssymb,amsthm,geometry,hyperref}
\geometry{margin=1in}
\newtheorem{theorem}{Theorem}
\newtheorem{proposition}[theorem]{Proposition}
\newtheorem{lemma}[theorem]{Lemma}
\newtheorem{corollary}[theorem]{Corollary}
\theoremstyle{definition}
\newtheorem{example}[theorem]{Example}
\newtheorem{remark}[theorem]{Remark}
\newcommand{\R}{\mathbb R}
\newcommand{\N}{\mathbb N}
\DeclareMathOperator{\NE}{NE}
\DeclareMathOperator{\cav}{cav}
\DeclareMathOperator{\supp}{supp}
\title{Grand-coalition arrival in the patient limit\\ GPT 6 Astra Ultra}
\author{Research attempt}
\date{9 October 2026}
\begin{document}
\maketitle

\section*{Target, source correction, and result}
The target asks whether a single fixed nested-coalition payoff model can have nonarriving equilibria for discount factors arbitrarily close to one. The source's revised equilibrium definition includes stays and compares a target with next-period continuation value. Its Remark 6.7 expressly withdraws the earlier deterministic-cycle example under that definition \cite{heitzig}. That example is not used here.

Two results are established: the bad discount factors in any fixed model have only finitely many interval components, and all equilibria in the two-player case arrive at the grand coalition. The first result changes a potentially sparse patient-limit counterexample search into a fixed-support interval search. It does not decide whether such an interval exists for any larger model.

\section*{Finite polynomial encoding of the equilibrium conditions}
Use the finite state set $X$ of laminar agreement families and the legal move sets from the target. A move $m$ at $x$ has target $\tau(m)$ and relevant set $R(m)\ne\varnothing$. All these sets are finite. Write $G\subseteq X$ for the grand-coalition states and $M(x)$ for the moving moves, with $M^+(x)$ also including all unilateral stays. Fix the real numbers $\pi_i(x)$ once and for all. For a stochastic matrix $p$ and $0<\delta<1$, define numbers $L_i(x),Z_i(x)$ by
\begin{equation}\label{bell}
 Z_i(x)=\sum_y p_{xy}L_i(y),\qquad
 L_i(x)=(1-\delta)\pi_i(x)+\delta Z_i(x).
\end{equation}
These equations have a unique solution: $(I-\delta p)^{-1}=\sum_{t\ge0}\delta^t p^t$ converges in the maximum row-sum norm. Thus $L_i$ is precisely the normalized discounted value and $Z_i$ the continuation value.

For targets $y,z$, define the strict comparison at $x$ by
\[
 y\succ_i^x z\ \Longleftrightarrow\
 [L_i(y)>L_i(z)]\ \lor\
 [L_i(y)=L_i(z),\ y=x,\ z\ne x],
\]
and define weak comparison as the negation of $z\succ_i^x y$. A move belongs to $C(x)$ exactly when
\[
 \bigwedge_{i\in R(m)}L_i(\tau(m))\ge Z_i(x)
\]
holds and there is no $m'\in M^+(x)$ with $R(m')\subseteq R(m)$ and $\tau(m')\succ_i^x\tau(m)$ for every $i\in R(m')$. A move in $C(x)$ belongs to $F(x)$ exactly when some player weakly prefers its target to the target of every member of $C(x)$.

Call $E(\delta,p,L,Z)$ the conjunction of stochasticity, \eqref{bell}, and the following conditions at every state:
\begin{align*}
 p_{xy}>0, y\ne x&\ \Longrightarrow\ \exists m\in C(x):\tau(m)=y,\\
 C(x)\cap M(x)\ne\varnothing&\ \Longrightarrow\ p_{xx}<1,\\
 m\in F(x)&\ \Longrightarrow\ p_{x,\tau(m)}>0.
\end{align*}
Every quantifier in the definitions of $C,F$ ranges over an explicitly finite set and can be expanded to a finite conjunction or disjunction. Therefore $E$ is a Boolean combination of polynomial equalities and strict or weak inequalities in $(\delta,p,L,Z)$, with the fixed payoffs as real coefficients. This also keeps the equality tie rule exact.

\begin{lemma}[Nonarrival certificate]
For a finite stochastic matrix $p$ and a nonempty target set $G$, some initial state has positive probability of never hitting $G$ if and only if there is a nonempty $C\subseteq X\setminus G$ such that
\[
 p_{xy}=0\quad(x\in C,\ y\notin C).
\]
\end{lemma}
\begin{proof}
If $X=G$, neither a nonarrival state nor a nonempty $C$ exists, so the equivalence holds. Otherwise, a chain starting in such a closed set $C$ never leaves it and therefore never hits $G$. Conversely, suppose no such set exists. In the directed support graph, every state outside $G$ has a path to $G$: otherwise the collection of states reachable from it would be a nonempty closed set outside $G$. Remove loops from each chosen path so its length is at most $|X|$. For each initial state outside $G$, the probability of following its path is positive. The minimum $\rho$ of these finitely many positive products is positive. Conditional on avoidance up to the beginning of any block of $|X|$ steps, the probability of hitting $G$ in that block is at least $\rho$. By the Markov property, avoidance for $k|X|$ steps has probability at most $(1-\rho)^k$, which tends to zero.
\end{proof}

\begin{theorem}[Fixed-model patient-limit dichotomy]
Fix any finite payoff model with the moves and equilibrium definition above, in particular any fixed model satisfying the target's surplus and sharing assumptions. Let
\[
 B=\{\delta\in(0,1):\text{some equilibrium has positive nonarrival probability from some state}\}.
\]
Then $B$ is a semialgebraic subset of the real line and is a finite union of points and intervals. Exactly one of the following alternatives holds:
\begin{enumerate}
\item There is $\delta_0<1$ with $B\cap(\delta_0,1)=\varnothing$.
\item There is $\delta_0<1$ with $(\delta_0,1)\subseteq B$. Moreover, one fixed nonempty set $C\subseteq X\setminus G$ and one fixed directed support $S\subseteq X\times X$ can be used for a witnessing equilibrium at every $\delta\in(\delta_0,1)$.
\end{enumerate}
For rational payoff data, it is decidable which alternative holds, with no polynomial-time bound asserted.
\end{theorem}
\begin{proof}
For each pair $(C,S)$, impose $E$, the closed-set equations of the lemma, and $p_{xy}>0$ for $(x,y)\in S$, $p_{xy}=0$ otherwise. Existentially quantify all coordinates of $p,L,Z$ to obtain a set $B_{C,S}$ of discounts.

We use the standard real quantifier-elimination theorem in exactly the following form: a set defined by a first-order formula with real polynomial comparisons is semialgebraic; for rational coefficients an equivalent quantifier-free formula can be computed. This is the Tarski--Seidenberg theorem and its effective form, stated in \cite{basu}, Section 2.1. This imported general theorem is not claimed as a new result here. It makes every $B_{C,S}$ semialgebraic, and the preceding lemma gives $B=\bigcup_{C,S}B_{C,S}$, a finite union.

A semialgebraic subset of the real line is a finite union of points and intervals: after quantifier elimination, gather the finitely many nonzero univariate polynomials in its formula. Each has finitely many roots. Between consecutive roots every polynomial has constant sign, so the formula has constant truth value. Zero polynomials do not create additional boundaries.

If $B$ has points arbitrarily close to one, finiteness of the union implies that some $B_{C,S}$ also has points arbitrarily close to one. Its finite interval decomposition must contain an interval ending at one; isolated points cannot accumulate there. This proves alternative 2 with fixed $C,S$. Otherwise alternative 1 holds. They cannot both hold. For rational data, apply effective quantifier elimination to the sentence
\[
 \exists d\ (0<d<1\ \wedge\ \forall\delta\,[d<\delta<1\Rightarrow\delta\in B]).
\]
Its truth selects alternative 2; its falsity selects alternative 1.
\end{proof}

This is a reduction for each fixed model, not a termination argument for searching over all player counts and all payoffs. Nor does fixed support impose a uniform positive lower bound on transition probabilities.

\section*{The two-player model}
\begin{proposition}[Arrival with two players at every discount]
If $N=\{1,2\}$ and the target's strict grand-surplus gap and positive sharing assumptions hold, every equilibrium reaches the grand coalition almost surely from every state, for every $0<\delta<1$.
\end{proposition}
\begin{proof}
There are only two states: $e=\varnothing$ and $g=\{N\}$. The assumptions give
$\pi_i(g)-\pi_i(e)=s_i(g,N)\gamma>0$ for both players. Put $a=p_{gg}$ and $b=p_{eg}$. Subtracting the two Bellman equations yields
\[
 L_i(g)-L_i(e)=
 \frac{(1-\delta)[\pi_i(g)-\pi_i(e)]}{1-\delta(a-b)}>0,
\]
because $a-b\in[-1,1]$ and the denominator is at least $1-\delta$.
At $g$, any termination to $e$ by player $i$ is dominated by that same player's stay at $g$. Thus no moving target is admissible at $g$, and equilibrium implies $p_{gg}=1$.
At $e$, merging to $g$ gives both relevant players at least $Z_i(e)$, since $Z_i(e)$ is a convex combination of $L_i(e)$ and $L_i(g)$. No stay can dominate this merge. It is the sole moving move: each player's stay at $e$ either fails profitability or is dominated only if allowed by the relevant-set rule, but in either case its target is strictly worse than $g$; hence the merge is a favorite whether or not stays remain in $C(e)$. The favorite condition forces $b>0$. Consequently the probability of remaining at $e$ for $t$ steps is $(1-b)^t\to0$.
\end{proof}

\section*{An adversarial limit check and unresolved gap}
The matrices $p_h=\left(\begin{smallmatrix}1-h&h\\0&1\end{smallmatrix}\right)$, $h>0$, all hit the second state almost surely, but their limit as $h\downarrow0$ is the identity and has nonarrival from the first state. This elementary calculation is a warning about a proposed compactness proof: arrival is not preserved when positive edges disappear. These matrices are only a Markov-chain check, not a claimed equilibrium counterexample.

The unresolved task is to find a fixed admissible model whose semialgebraic bad set contains a final interval, or to exclude every such interval for every player count. No rational model was run through full quantifier elimination in this attempt.

\begin{thebibliography}{9}
\bibitem{heitzig} J. Heitzig, \emph{Does the grand coalition form? Persistence, arrival, and the role of the sharing rule in a dynamic process of nested binding agreements}, arXiv:2608.17766v2 (2026), Definitions 2.14--2.15, Remark 6.7, and Section 8. \url{https://arxiv.org/html/2608.17766v2}. These revised definitions and the explicit correction to the old cycle were checked on 9 October 2026.
\bibitem{basu} S. Basu, \emph{Algorithms in Real Algebraic Geometry: A Survey}, arXiv:1409.1534 (2014), Section 2.1, Theorem 2.1. \url{https://arxiv.org/pdf/1409.1534}. Used solely for real quantifier elimination; all equilibrium-specific encoding and finite-chain arguments are proved above.
\end{thebibliography}

\end{document}
